\chapter{Robust Hypothesis Testing}
\label{sec:p15-robust-hypothesis-testing}

A nominal error probability concerns a specified sampling law.
Robust testing instead specifies a class of possible laws and controls
the largest error over that class. Huber's original construction clips
the individual likelihood ratios and obtains an exact least-favorable
pair for two particular kinds of neighborhoods
\cite[Secs.~2--4 and 7]{Huber1965RobustProbabilityRatio}.
The neighborhoods remain classes of probability measures, including
singular contaminations; they are not replaced by a finite list of models.

\section{The experiment and its two uncertainty classes}
\label{sec:p15-robust-neighborhoods}

\begin{hbdefinition}{def:p15-robust-test-experiment}{Worst-case errors for a fixed sample size}
Let \(P_0\ne P_1\) be probabilities on a nonempty measurable space
\((\mathsf X,\mathscr X)\). For nonempty classes \(\mathcal C_0,\mathcal C_1\),
observe an integer \(n\geq1\) independent observations with one common law
\(R\in\mathcal C_i\) under hypothesis \(i\).
A test is measurable \(\varphi:\mathsf X^n\to[0,1]\), with rejection
probability \(E_{R^{\otimes n}}\varphi\).
Define
\[
 A(\varphi)=\sup_{R\in\mathcal C_0}E_{R^{\otimes n}}\varphi,\qquad
 B(\varphi)=\sup_{R\in\mathcal C_1}E_{R^{\otimes n}}(1-\varphi).
\]
Its worst-case power is \(1-B(\varphi)\). The level-\(\alpha\) problem
minimizes \(B\) subject to \(A\leq\alpha\).
For losses \(L_0,L_1>0\), Huber also considers minimizing
\(\max\{L_0A,L_1B\}\), or the weighted risk
\(w_0A+w_1B\), \(w_i\geq0\), not both zero.
The latter includes \(w_i=\lambda_iL_i\) for specified hypothesis
prior probabilities \(\lambda_i\geq0\), \(\lambda_0+\lambda_1=1\).
\end{hbdefinition}

\begin{hbdefinition}{def:p15-huber-neighborhoods}{Contamination and the source's total-variation convention}
For \(0\leq\varepsilon_i<1\), the contamination classes are
\[
 \mathcal C_i=
 \{(1-\varepsilon_i)P_i+\varepsilon_i H:H
          \text{ is any probability on }(\mathsf X,\mathscr X)\}.
\]
Write \(a_i=1-\varepsilon_i\).
For probabilities \(P,R\), define
\[
 d_{\rm TV}(P,R)=\sup_{E\in\mathscr X}|P(E)-R(E)|,\qquad
 \|P-R\|_{\rm var}=2d_{\rm TV}(P,R).
\]
Huber's second construction uses the equal-radius classes
\[
 \mathcal V_i=\{R:\|R-P_i\|_{\rm var}\leq\varepsilon\},
                 \qquad 0\leq\varepsilon<2.
\]
Thus their radius in the sup-event convention is \(\delta=\varepsilon/2\).
The two uses of \(\varepsilon\) describe different neighborhoods.
For a real number \(z\), write \(z_+=\max(z,0)\).
\end{hbdefinition}

\begin{hbtheorem}{thm:p15-robust-overlap}{proved=annex:p15.robust.overlap}{Overlap and the exact loss of uniform discrimination}
Let \(\mu=P_0+P_1\), choose finite densities \(p_i=dP_i/d\mu\), and put
\[
 m=\int\max(a_0p_0,a_1p_1)\,d\mu,\qquad d=d_{\rm TV}(P_0,P_1)>0.
\]
The contamination classes intersect exactly when \(m\leq1\).
For equal contamination fractions \(\varepsilon_0=\varepsilon_1=e\),
they are disjoint exactly when \(e<d/(1+d)\).
The total-variation classes intersect exactly when \(\varepsilon\geq d\).

Whenever the null and alternative classes intersect, the optimal
worst-case power at level \(\alpha\in[0,1]\) is \(\alpha\), attained
by the constant randomized test \(\varphi=\alpha\), for every \(n\).
The minimax maximum loss is \(L_0L_1/(L_0+L_1)\), attained by
\(\varphi=L_1/(L_0+L_1)\). The minimum weighted worst-case risk is
\(\min(w_0,w_1)\), attained by an appropriate constant decision.
\end{hbtheorem}

\section{Normalization and clipping for contamination}
\label{sec:p15-contamination-clipping}

Use \(\mu=P_0+P_1\). Define the nominal ratio \(r=p_1/p_0\) where
\(p_0>0\), let \(r=\infty\) where \(p_0=0<p_1\), and set \(r=1\)
where both densities vanish. The last choice matters to arbitrary
contaminations even though it has no effect under the nominal laws.

\begin{hbtheorem}{thm:p15-huber-contamination-pair}{proved=annex:p15.huber.contamination.pair}{The original normalization equations and clipped pair}
Suppose \(m>1\). Define
\[
 F(u)=\int\max(p_1,up_0)\,d\mu,\qquad
 G(v)=\int\max(p_0,vp_1)\,d\mu,\quad u,v\geq0.
\]
Choose \(u,v\) by
\[
 F(u)=a_1^{-1},\qquad G(v)=a_0^{-1},
\]
using \(u=0\) when \(\varepsilon_1=0\), and \(v=0\) when
\(\varepsilon_0=0\). When its contamination fraction is positive,
each equation has a unique positive solution.
Set
\[
 q_0=a_0\max(p_0,vp_1),\qquad
 q_1=a_1\max(p_1,up_0),\qquad b=a_1/a_0.
\]
These are probability densities of \(Q_i\in\mathcal C_i\).
Moreover \(bu<1\), \(v<b\), hence \(u<1/v\), with \(1/0=\infty\).
Put \(c'=u,c''=1/v\). The score
\[
 \ell(x)=b\max\{c',\min(c'',r(x))\}
\]
equals \(q_1/q_0\) on the support of the pair, with the usual extended
ratio convention. It is assigned the displayed clipped value on the
common zero-density set as well.
These are Huber's equations (2)--(4).

If both contamination fractions are positive, then
\(0<bc'<1<bc''<\infty\). If only \(\varepsilon_0=0\), the upper endpoint
may be infinite but the lower endpoint is positive.
If only \(\varepsilon_1=0\), the lower endpoint is zero but the upper
endpoint is finite. If both are zero, the problem is the ordinary
two-law test, including its singular-support cases.
\end{hbtheorem}

\begin{hbtheorem}{thm:p15-huber-tv-pair}{proved=annex:p15.huber.tv.pair}{The distinct total-variation construction}
Suppose \(0<\varepsilon<d\), and write \(s=p_0+p_1\).
There are unique \(k_-,k_+\in[0,1]\) satisfying
\[
 \int(k_-s-p_1)_+\,d\mu=\varepsilon/2,\qquad
 \int(p_1-k_+s)_+\,d\mu=\varepsilon/2.
\]
They obey \(0<k_-<1/2<k_+<1\). Define
\[
 q_1=\max\{k_-s,\min(k_+s,p_1)\},\qquad q_0=s-q_1.
\]
Then \(Q_i\) are probabilities and
\(\|Q_i-P_i\|_{\rm var}=\varepsilon\).
With \(c'=k_-/(1-k_-)\), \(c''=k_+/(1-k_+)\),
\[
 \ell=\max\{c',\min(c'',r)\}=q_1/q_0
\]
on the pair's support. On its common null set use the displayed score,
which lies between the clipping endpoints.
Equivalently, on either clipped region the source's densities are
\(q_0=s/(1+c)\), \(q_1=cs/(1+c)\), using the relevant endpoint \(c\);
on the middle region \(q_i=p_i\).
At \(\varepsilon=0\) use \(Q_i=P_i\) and the nominal test.
For \(\varepsilon\geq d\) use the overlap conclusion rather than
crossing the clipping endpoints.
\end{hbtheorem}

\section{The least-favorable order and the minimax proof}
\label{sec:p15-robust-minimax}

\begin{hblemma}{thm:p15-huber-stochastic-order}{proved=annex:p15.huber.stochastic.order}{Huber's order for every allowed contamination}
For either constructed separated pair, let \(R_i\) range over its
corresponding neighborhood. For every real \(t\),
\[
 R_0\{\ell<t\}\geq Q_0\{\ell<t\}
       \geq Q_1\{\ell<t\}\geq R_1\{\ell<t\}.
\]
These inequalities include singular contaminations on nominal null sets.
When a score endpoint is infinite, the order is understood on
\([0,\infty]\); equivalent weak-threshold inequalities follow by limits.
\end{hblemma}

\begin{hblemma}{thm:p15-robust-product-order}{proved=annex:p15.robust.product.order}{A complete real-score coupling and product transfer}
For nonnegative extended scores whose distributions have the preceding
order, independent copies can be coupled so that, coordinate by coordinate,
the null-neighborhood scores are no larger than the \(Q_0\) scores,
and the alternative-neighborhood scores are no smaller than the \(Q_1\)
scores. This conclusion requires no regular conditional law on
\(\mathsf X\); the coupling is constructed on the real score space.

If at least one contamination fraction is positive, or the total-variation
radius is positive, the product \(T_n=\prod_{j=1}^n\ell(X_j)\) is
well defined in \([0,\infty]\): zero and infinity cannot both occur
among its possible factors. For \(c\in[0,\infty)\), \(0\leq\kappa\leq1\),
let
\[
 \varphi_{c,\kappa}
   ={\bf1}_{\{T_n>c\}}+\kappa{\bf1}_{\{T_n=c\}} .
\]
Then
\[
 \sup_{R\in\mathcal C_0}E_{R^{\otimes n}}\varphi_{c,\kappa}
       =E_{Q_0^{\otimes n}}\varphi_{c,\kappa},\qquad
 \inf_{R\in\mathcal C_1}E_{R^{\otimes n}}\varphi_{c,\kappa}
       =E_{Q_1^{\otimes n}}\varphi_{c,\kappa},
\]
with \(\mathcal V_i\) in place of \(\mathcal C_i\) in the variation case.
The boundary randomization is constant on the entire threshold set.
Arbitrary data-dependent changes within that set are not covered by
this monotonicity argument.
\end{hblemma}

\begin{hbtheorem}{thm:p15-huber-minimax-tests}{proved=annex:p15.huber.minimax.tests}{All three fixed-sample minimax decisions}
For either separated neighborhood pair, the following tests solve the
three problems in \cref{def:p15-robust-test-experiment}.
\begin{enumerate}
\item At \(0<\alpha<1\), choose a threshold and constant boundary
randomization with \(E_{Q_0^{\otimes n}}\varphi=\alpha\).
The resulting likelihood-ratio test maximizes worst-case power subject
to \(A\leq\alpha\), and its two worst-case errors occur at \(Q_0,Q_1\).
At level zero use \({\bf1}_{\{T_n=\infty\}}\), and at level one use 1.
\item For \(w_0,w_1>0\), reject when \(T_n>w_0/w_1\), with any constant
boundary randomization. This minimizes \(w_0A+w_1B\). If a weight
vanishes, the corresponding constant decision is optimal.
\item For \(L_0,L_1>0\), choose a likelihood threshold and constant
boundary randomization satisfying
\[
 L_0E_{Q_0^{\otimes n}}\varphi
       =L_1E_{Q_1^{\otimes n}}(1-\varphi).
\]
Such a choice exists, including the possible zero-risk singular case,
and minimizes \(\max\{L_0A,L_1B\}\).
\end{enumerate}
For zero radii in both neighborhoods, use the ordinary product-density
test from \cref{thm:p15-neyman-pearson}. On sample tuples containing
both a nominal-only-null and a nominal-only-alternative coordinate,
both product densities are zero; give the test any fixed value there
instead of multiplying zero by infinity.
\end{hbtheorem}

For unequal contamination fractions,
\[
 \log T_n=n\log b+\sum_{j=1}^n
     \log\max\{c',\min(c'',r(X_j))\}.
\]
For fixed \(n\), the factor \(b^n\) can be absorbed into the decision
threshold. In a sequential test \(n\) changes, so this deterministic
per-observation drift must be retained.

\section{Neighborhood sensitivity and complete calculations}
\label{sec:p15-robust-examples}

\begin{hbtheorem}{thm:p15-robust-sensitivity}{proved=annex:p15.robust.sensitivity}{Larger neighborhoods and the source's sign-statistic limit}
For one observation and any measurable \(\varphi:\mathsf X\to[0,1]\),
\[
 A(\varphi)=a_0E_{P_0}\varphi+\varepsilon_0\sup_{x\in\mathsf X}\varphi(x),
 \quad
 1-B(\varphi)=a_1E_{P_1}\varphi+\varepsilon_1\inf_{x\in\mathsf X}\varphi(x).
\]
The extrema are pointwise, not essential extrema under a nominal law.
An unattained supremum is still the supremum over contaminations.

Increasing either contamination fraction, or the variation radius,
enlarges the respective neighborhood. The optimal worst-case power
at a fixed level cannot increase when either class is enlarged.
For contamination, \(c'\) is nondecreasing in \(\varepsilon_1\) and
\(c''\) is nonincreasing in \(\varepsilon_0\), for the normalized
choices while the classes remain separated.

For equal positive contamination \(e<e_*=d/(1+d)\),
\(c'\uparrow1\), \(c''\downarrow1\) as \(e\uparrow e_*\).
For every \(x\) with \(r(x)\ne1\),
\[
 \frac{\log\max(c',\min(c'',r(x)))-\log c'}
      {\log c''-\log c'}\longrightarrow {\bf1}_{\{r(x)>1\}}.
\]
For a fixed finite sample avoiding \(r=1\), the sum therefore tends
to its count of ratios exceeding one. This holds almost surely under
\(P_i^{\otimes n}\) if \(P_i(r=1)=0\).
It is a limit of the rescaled statistic; at the overlap boundary the
maximin power is still only the test's level.
\end{hbtheorem}

\begin{hbtheorem}{thm:p15-robust-testing-calculation}{proved=annex:p15.robust.testing.calculation}{Finite and normal contamination calculations}
On \(\{a,b,c\}\), let
\[
 P_0=(3/5,3/10,1/10),\qquad P_1=(1/10,3/10,3/5).
\]
With contamination fractions \(1/10\), the least-favorable laws are
\[
 Q_0=(54,27,19)/100,\qquad Q_1=(19,27,54)/100,
 \qquad (c',c'')=(19/54,54/19).
\]
For one observation at level \(1/10\), reject only at \(c\), with
probability \(10/19\). Its worst-case size is \(1/10\), and its
worst-case power is \(27/95\).
The nominal test that always rejects at \(c\) has nominal size \(1/10\)
but worst-case contaminated size \(19/100\).

For two independent observations, the product score and its law under
the least-favorable pair are
\[
\begin{array}{c|ccccc}
t&(19/54)^2&19/54&1&54/19&(54/19)^2\\ \hline
10000\,Q_0^{\otimes2}(T_2=t)&2916&2916&2781&1026&361\\
10000\,Q_1^{\otimes2}(T_2=t)&361&1026&2781&2916&2916
\end{array}
\]
At level \(1/10\), reject at the largest score and randomize with
probability \(71/114\) at score \(54/19\). The worst-case power is
\(8991/19000\).

For equal contamination \(0<e<1/3\), the same example has
\[
 Q_0=(.6(1-e),.3(1-e),.1+.9e),\quad Q_1
       \text{ with first and last coordinates exchanged},
 \quad c'=\frac{1+9e}{6(1-e)},\quad c''=1/c'.
\]
At \(e=1/3\) the classes meet.
For the different total-variation neighborhoods with
\(\varepsilon=1/5\), equivalently sup-event radius \(1/10\), the pair is
\[
 Q_0=(1/2,3/10,1/5),\qquad Q_1=(1/5,3/10,1/2),
 \qquad(c',c'')=(2/5,5/2).
\]
Its one-observation level-\(1/10\) robust test rejects at \(c\) with
probability \(1/2\) and has worst-case power \(1/4\).

For a continuous example let \(P_0=\operatorname{normal}(-1,1)\),
\(P_1=\operatorname{normal}(1,1)\), and let \(\Phi\) be the standard
normal distribution function. Put
\[
 D=\Phi(2)+\tfrac12\exp(-2),\qquad \varepsilon_{\rm N}=1-D^{-1}.
\]
These are separated equal-contamination neighborhoods with
\(0<\varepsilon_{\rm N}<1\).
Their clipping endpoints are \(\exp(-2),\exp(2)\), corresponding to clipping
the observation at \(-1,1\). The least-favorable null Lebesgue density is
\[
 q_0(x)=D^{-1}
 \begin{cases}
 (2\pi)^{-1/2}\exp(-(x+1)^2/2),&x<1,\\
 \exp(-2)(2\pi)^{-1/2}\exp(-(x-1)^2/2),&x\geq1,
 \end{cases}
 \qquad q_1(x)=q_0(-x).
\]
Rejecting at \(x>0\) equalizes the two worst-case errors at
\(1-\Phi(1)/D\), so is minimax for equal losses.
Its nominal null error is only \(1-\Phi(1)\).
\end{hbtheorem}

\begin{hbtheorem}{thm:p15-robust-sampling-boundary}{proved=annex:p15.robust.sampling.boundary}{Why the sample-level neighborhood matters}
The product neighborhood
\(\{R^{\otimes n}:R=(1-e)P+eH\}\) and the whole-sample neighborhood
\(\{(1-e)P^{\otimes n}+eK:K\text{ any joint sample law}\}\) are different.
For \(P=\delta_0,H=\delta_1,n=2\), independent per-observation
contamination gives probability \(2e(1-e)>0\) of unequal coordinates
when \(0<e<1\). The whole-sample law with \(K=\delta_{(1,1)}\)
assigns that event probability zero.
Moreover arbitrary dependence with correct contaminated marginals
is not licensed by the product theorem: in the three-value example,
two perfectly repeated draws from \(Q_0\) give rejection probability
\(19/100\) for the displayed two-observation robust test, exceeding
its product-sample level \(1/10\).
\end{hbtheorem}

\section{The source's extensions and their exact premises}
\label{sec:p15-robust-source-extensions}

\begin{hbtheorem}{thm:p15-robust-mlr-extension}{proved=annex:p15.robust.mlr.extension}{Separated ordered composite hypotheses}
Let \(P_\theta\), \(\theta\) in an ordered parameter set, have finite
Borel densities \(p_\theta\) on \(\mathbb R\) against a common
\(\sigma\)-finite Borel measure \(\nu\). Assume the pointwise
monotone-likelihood-ratio condition
\[
 p_{\theta'}(y)p_\theta(x)\geq
 p_{\theta'}(x)p_\theta(y)
       \quad(x<y,\ \theta<\theta').
\]
Fix \(\theta_0<\theta_1\) and separated endpoint neighborhoods as above.
Extend the endpoint nominal ratio monotonically across its common
zero-density set, and clip it as in the relevant construction.
Then the same least-favorable pair and fixed-sample minimax tests apply
to the unions of the neighborhoods about all \(P_\theta\) with
\(\theta\leq\theta_0\), and all with \(\theta\geq\theta_1\).
The annex constructs the monotone extension and proves the needed
stochastic-order implication, allowing varying supports.
The sequential results in
\cref{thm:p15-robust-sequential-errors,thm:p15-robust-sequential-asymptotic}
also apply to these union classes under their stated premises.
\end{hbtheorem}

\begin{hbtheorem}{thm:p15-robust-sequential-errors}{proved=annex:p15.robust.sequential.errors}{Sequential errors and finite expected termination}
For a separated constructed pair \(Q_0\ne Q_1\), write
\(\gamma=\log\ell\), with \(\log0=-\infty,\log\infty=\infty\).
For constants \(A<0<B\), stop at
\[
 N=\inf\{n\geq1:\gamma(X_1)+\cdots+\gamma(X_n)\notin(A,B)\}.
\]
Reject the null at the upper exit and accept it at the lower exit.
An infinite increment causes immediate exit; no later increment is
added to it. Put \(\eta=\int\sqrt{q_0q_1}\,d\mu<1\).
For every i.i.d. neighborhood law the stopping time is almost surely
finite with finite expectation. In fact, for \(n\geq1\),
\[
 \Pr_{R_0}(N>n)\leq e^{-A/2}\eta^n,\qquad
 \Pr_{R_1}(N>n)\leq e^{B/2}\eta^n .
\]
The upper-exit error is largest at
\(Q_0\), and the lower-exit error is largest at \(Q_1\).
These statements retain the factor \(b\) in the contamination score.
They compare error probabilities; an assertion of finite-boundary
minimax expected sample size is a different claim.
\end{hbtheorem}

\begin{hbtheorem}{thm:p15-robust-sequential-asymptotic}{proved=annex:p15.robust.sequential.asymptotic}{The expected-size asymptotic under bounded clipping}
Assume in addition \(0<\ell_-\leq\ell\leq\ell_+<\infty\),
and let \(A\to-\infty\),
\(B\to\infty\), with no relative growth restriction. For each null
neighborhood law \(R_0\) and alternative law \(R_1\),
\[
 E_{R_0}N\sim\frac{|A|}{|E_{R_0}\gamma|},\qquad
 E_{R_1}N\sim\frac{B}{E_{R_1}\gamma}.
\]
The relative errors tend to zero uniformly over their respective
neighborhoods. In particular
\[
 \frac{\sup_{R_0}E_{R_0}N}{E_{Q_0}N}\longrightarrow1,\qquad
 \frac{\sup_{R_1}E_{R_1}N}{E_{Q_1}N}\longrightarrow1.
\]
The proof supplies the stopping identity and overshoot bounds needed
for the source's expected-size argument, rather than assuming them.
\end{hbtheorem}

Huber's paper lists questions about finite-boundary expected sample
size, a shared unknown contamination law, and more than two decisions
\cite[Sec.~8]{Huber1965RobustProbabilityRatio}. They are not settled by
the fixed-sample minimax theorem or the displayed asymptotic.
The unbounded and singular endpoints are completed in
\cref{thm:p15-robust-endpoint-closure}. That result proves the full
least-favorable expected-size ratio under vanishing errors.
Its finite-entropy, infinite-entropy and singular-mass size classifications
apply when both boundaries diverge. The corresponding leading minimax
cost distinctions are in \cref{thm:p15-sequential-leading-minimax}.
Robust estimation, influence functions, breakdown, and inference after
choosing a neighborhood radius from the data are further developments.
